\documentclass[../../main.tex]{subfiles}
\begin{document}

\begin{center}
  {\fontsize{48pt}{58pt}\selectfont \textbf{1976}\normalsize\textbf{年度}}
\end{center}

\begin{table}[htb]
  \centering
  \caption{出題テーマと難易度}
  \begin{tabular}{|c|p{9cm}|c|}
    \hline
    \textbf{問題番号} & \textbf{テーマ} & \textbf{難易度} \\
    \hline
    第1問 & 整式の剰余 & \\
    \hline
    第2問 & 不等式の恒等成立条件 & \\
    \hline
    第3問 & 図形の軌跡と面積 & \\
    \hline
    第4問 & 積分方程式 & \\
    \hline
    第5問 & ベクトルの内積 & \\
    \hline
    第6問 & 一次変換 & \\
    \hline
  \end{tabular}
\end{table}

\begin{table}[htb]
  \centering
  \caption{解答}
  \renewcommand{\arraystretch}{1.4}
  \begin{tabular}{|c|c|p{9cm}|}
    \hline
    \textbf{問題番号} & \textbf{小問} & \textbf{答え} \\
    \hline
    第1問 & - & $r(x)=x(x-1)$ \\
    \hline
    第2問 & $0<c\le\dfrac{1}{2}$ & $b>c$ \\
    \cline{2-3}
    & $\dfrac{1}{2}\le c$ & $b>c^2+\dfrac{1}{4}$ \\
    \hline
    第3問 & - & $S=\dfrac{a^2}{3}(4\sqrt{2}-5)$ \\
    \hline
    第4問 & - & $f(x)=e^x-2e^{2x}$（$a=\dfrac{1}{3-2e}$） \\
    \hline
    第5問 & (1) & $|\vec a|=1,\ |\vec b|=\sqrt{2},\ |\vec c|=\sqrt{3}$ \\
    \cline{2-3}
    & (2) & $\vec a\cdot\vec b=\vec a\cdot\vec c=\vec b\cdot\vec c=0$ \\
    \hline
    第6問 & - &  \\
    \hline
  \end{tabular}
\end{table}

\paragraph{試験形態}
6問．試験時間は要確認．

\end{document}
